\chapter{Meta-conhecimento}
\label{cap:metaconhecimento}

\section{Definição e finalidade}
\label{sec:meta-definicao}

Neste projeto, meta-conhecimento é o conhecimento sobre as próprias regras, e não
sobre a condução em si, usado para controlar e acelerar a inferência. Saber que
se deve reduzir a velocidade quando o limite é ultrapassado é conhecimento do
domínio, enquanto saber que a regra que o faz lê os factos \cod{ego_speed/1} e
\cod{speed_limit/1}, que tem precedência 300 e que já disparou neste ciclo é
meta-conhecimento.

Serve três fins: a resolução de conflitos, a terminação (refração) e a
eficiência (índice de dependências).

\section{Representação no motor Prolog}
\label{sec:meta-representacao}

O meta-conhecimento está em predicados dinâmicos de \cod{engine.pl}, declarados com \cod{:- dynamic}~\cite{swideclare} e separados da
memória de trabalho (\cod{wm_fact/2}, \cod{wm_support/3}, etc.), cujos
predicados têm o prefixo \cod{wm_} (\cref{tab:metaconhecimento}).

\begin{table}[H]
  \centering
  \caption{Predicados de meta-conhecimento do motor}
  \label{tab:metaconhecimento}
  \small
  % Colunas de texto em bandeira (sem justificação), para não abrir espaços.
  \begin{tabularx}{\textwidth}{@{}L{4.7cm}VL{3.9cm}V>{\raggedright\arraybackslash}X@{}}
    \toprule
    \textbf{Predicado} & \textbf{Conhecimento representado} & \textbf{Utilização} \\
    \midrule
    \cod{rule_reads(Nome/Aridade, Regra)} & Tipos de facto que cada regra lê
      & Avaliar de novo só as regras afetadas (\cod{fire/2}, \cod{rematch/1}) \\ \linhasep
    \cod{rule_precedence(Regra, P)} & Precedência de cada regra
      & Escolher a ativação seguinte (\cod{best_activation/2}) \\ \linhasep
    \cod{rule_rank(Regra, N)} & Posição da regra na base de conhecimento
      & Desempate determinístico entre regras com a mesma precedência \\ \linhasep
    \texttt{fired(Regra, Condições)} & Instanciações que já dispararam no ciclo
      & Refração (\cod{already_fired/2}) \\
    \bottomrule
  \end{tabularx}
\end{table}

Este meta-conhecimento é gerado dinamicamente, e não escrito manualmente, em dois sentidos:

\begin{enumerate}
  \item \textbf{É derivado das regras durante a execução.} Quando o motor
    arranca, \cod{ensure_index/0} calcula \cod{rule_reads/2},
    \cod{rule_precedence/2} e \cod{rule_rank/2} a partir das próprias regras
    (\cref{sec:meta-indice}) e guarda-os com \cod{assertz}~\cite{swidynamic}.
    Uma regra acrescentada ou alterada é refletida no arranque seguinte, sem
    manutenção manual, e o mesmo passo verifica que não há identificadores
    repetidos.
  \item \textbf{Muda durante a inferência.} \cod{fired/2} ganha uma entrada a
    cada disparo e é apagado no início de cada ciclo.
\end{enumerate}

\section{Índice de dependências}
\label{sec:meta-indice}

O índice (\cod{rule_reads/2}) regista, para cada regra, os tipos de facto
(\cod{Nome/Aridade}) de que depende. \cod{ensure_index/0} obtém-nos por análise
estática da regra:

\begin{itemize}
  \item cada padrão e cada padrão negado (\cod{not/1}) contam com o seu
    \cod{Nome/Aridade};
  \item em \cod{test/1} e \cod{prove/1}, contam os factos lidos com
    \cod{known/1}, no próprio objetivo e nas cláusulas dos predicados auxiliares
    que ele chama, seguidos transitivamente com \cod{clause/2}.
\end{itemize}

Com a base de conhecimento atual, o índice tem $\sim$600 entradas, para $\sim$120
regras e $\sim$66 tipos de facto. A regra
\cod{"Avoid pothole in centre lane by changing right when right is clear"}, por
exemplo, depende de nove tipos de facto, mas só três aparecem nas suas condições
(\cod{ego_lane/1}, \cod{pothole/4} e \cod{ego_speed/1}). Os outros seis, como as
marcações da estrada, a ocupação das vias e o estado da manobra, vêm dos
predicados auxiliares, sobretudo de \cod{safe_to_change_lane/1}, chamado através
de \cod{prove/1}. Um índice feito só com os padrões visíveis da regra não veria
estas dependências.

Durante a inferência, \cod{fire/2} recolhe os tipos dos factos novos que o
disparo acrescentou e, através de \cod{rule_reads/2}, avalia de novo apenas as
regras que os leem (\cod{rematch/1}). Sem o índice, todas as
regras teriam de ser avaliadas de novo depois de cada disparo.

O teste \cod{DependencyIndexBenchmarkTest} mede o efeito do índice. Corre os 14
cenários da biblioteca durante 150 ciclos cada (2\,100 ciclos) em dois processos
Prolog lado a lado, um com o índice e outro sem ele. Em todos os ciclos, as duas
execuções tomaram a mesma decisão e dispararam as mesmas regras (10\,231
disparos): o índice só muda o custo, que o próprio motor mede
(\cod{cycle_stats/3}). Com o índice, a inferência ocupa em média cerca de
1,5\,\% do ciclo de 100\,ms (\cref{tab:desempenho}). Os tempos cobrem só a
inferência, sem a comunicação JSON, e foram medidos numa única máquina (Intel
Core Ultra 7 258V, SWI-Prolog 9.2.9).

\begin{table}[H]
  \centering
  \caption{Efeito do índice de dependências (14 cenários, 2\,100 ciclos)}
  \label{tab:desempenho}
  \small
  \begin{tabularx}{\textwidth}{@{}XVC{3.1cm}VC{3.1cm}VC{1.9cm}@{}}
    \toprule
    & \textbf{Com índice} & \textbf{Sem índice} & \textbf{Redução} \\
    \midrule
    Avaliações de regras, total (por ciclo) & 258\,894 (123) & 1\,479\,720 (705)
      & 82,5\,\% \\ \linhasep
    Tempo médio por ciclo & 1,47\,ms & 6,41\,ms & 77,0\,\% \\ \linhasep
    Pior cenário (S10, chuva) & 2,64\,ms & 9,62\,ms & 72,6\,\% \\
    \bottomrule
  \end{tabularx}
\end{table}

\section{Conhecimento de controlo}
\label{sec:meta-controlo}

Entre as ativações prontas, \cod{best_activation/2} escolhe a de maior
precedência (\cod{rule_precedence/2}). Em empate, escolhe a da regra declarada
primeiro (\cod{rule_rank/2}) e, depois, a correspondência encontrada primeiro. A
precedência de cada regra é a \emph{salience} da regra Drools que ela replica,
para que os dois motores do sistema disparem pela mesma ordem. As únicas
exceções são as duas regras intermédias de peões, que só existem em Prolog
(\cref{lst:regras-ou}).

A refração usa \cod{fired/2}. Antes de uma instanciação entrar na agenda,
\cod{already_fired/2} compara as suas condições com as das instanciações que já
dispararam, por variante (\cod{=@=}), e deixa de fora as repetidas. Como os
factos nunca são retirados dentro de um ciclo, isto impede que a mesma conclusão
seja produzida repetidamente. Um limite de 5000 disparos por ciclo
(\cod{max_firings_per_cycle/1}) protege ainda contra uma base de regras mal construída,
que não pare de gerar instanciações novas, devolvendo um erro em vez de bloquear a
simulação.
